% Section 7 -- Case study
\section{Case study: the Orange Spain credential hijack, 2024-01-03}
\label{sec:case}

The case study is chosen because it is the event where the gap between
what happened and what the telemetry records is widest, and because it
shows how the system behaves when the control plane moves and the data
plane does not.

\subsection{What happened}

On 2024-01-03 an attacker used stolen RIPE NCC credentials---the
account had no second factor---to log into Orange Spain's (AS12479)
RIPE account and publish invalid RPKI ROAs for Orange Spain address
space, pointing the origins at AS49581, which is not an Orange
ASN~\cite{apnic-orange-2024}. Networks performing RPKI origin
validation began rejecting Orange Spain announcements at approximately
14:20 UTC; AS5511, Orange Spain's transit provider, propagated the
affected /12s. Reachability was mostly restored by 18:00 UTC, after the
invalid ROAs were removed. Roughly two million addresses were affected.

Three features make this event a good test. The cause is
\emph{control-plane initiated and data-plane visible}: a routing
configuration change produced a reachability effect. The victim is the
announcer, so the detector's ``who is the attacker'' question and the
ground truth's ``who was harmed'' question point at the same AS. And the
ground truth records \emph{no traffic anomaly} for the window
(\texttt{has\_traffic\_signal = false}, \texttt{traffic\_hiz = 0}):
the outage was visible to validating networks as unreachability, not to
the vantage point's traffic series as a volume change.

\subsection{What the cache contains}

We replay the cached bundle for this event (window 14:45--18:00 UTC,
the closest cached window to the incident).

\textbf{Control plane.} The routing block holds 9{,}158
announcement-level detection records spanning 14:54:29 to 18:04:20, all
typed \texttt{origin\_hijacking}. The merged timeline that the engine
reads is a bounded view of that list: it retains the first $50$ entries,
of which $49$ are routing records---all stamped 14:54:29---and one is
the synthesized traffic record discussed below. The retained count is not
a measurement of hijack intensity. Every record carries
\texttt{expected\_origin: null}---the window was collected in the
pipeline's light mode, in which the expected-origin baseline is not
built---so the detector flags \emph{the victim's own announcements} and
attributes them to \texttt{hijacker\_as = 12479} with
\texttt{victim\_as = null} and \texttt{role = attacker}. The cache's own
summary fields disagree with its lists in the same way: it reports
\texttt{total\_prefix\_hijacks = 0} and
\texttt{leak\_analysis\_skipped = true} with the reason
\texttt{light\_mode}, while the hijack list it contains is non-empty. The
ground truth, by contrast, records five hijacked events with
\texttt{hijacker\_as = AS5511} between 14:39 and 16:03 and 220 routing
events in the window. The cache and the ground truth therefore disagree
in both magnitude and role: 9{,}158 flags against the victim's own AS,
versus 5 events against the transit provider.

\textbf{Data plane.} The traffic block holds 134 five-minute data
points, \texttt{anomaly\_count = 0}, and a baseline that was never
built: \texttt{baseline\_weeks\_available} = 0,
\texttt{baseline\_method} = \texttt{non\_periodic\_fallback}, periodicity
confidence 0.0, and \texttt{historical\_avg = 0.0}. The same block
reports a peak-to-trough drop of 81.7\,\% over the expanded window. The
``no anomaly'' verdict is therefore a statement about the baseline, not
about the traffic---a distinction the cache preserves and the aggregate
count does not. The single traffic record in the correlation timeline is
timestamped 07:22:30, outside the incident window, and carries
\texttt{cluster\_anomaly\_count = 0} and \texttt{pct\_change = 0.0}
while reporting severity 0.82 and
\texttt{synthesized\_from = expanded\_window}: it is a window-midpoint
synthesis rather than an observation.

\begin{figure}[t]
\centering
\includegraphics[width=0.9\linewidth]{figs/fig_case_orange.pdf}
\caption{The cached window for the Orange Spain event on five-minute
buckets. (a) The routing block as recorded---all 9{,}158
announcement-level detections, 25 of the 39 buckets non-empty---where
the density is created by a missing expected-origin baseline in light
mode, each record flagging the victim's own announcement, rather than
by the scale of the hijack. The merged timeline the engine reads is a
bounded view of this block that retains the first $50$ entries, all
stamped 14:54:29 (\S\ref{sec:eval-shape}), so the routing side collapses
to a single spike. (b) What Phase~2 sees of the same window on its own
five-minute grid: one routing point at bucket 90 and the synthesized
traffic record at bucket 0, with the other 89 buckets empty and floored
at $\varepsilon$. Neither series has an interior, so the control plane
has no sustained deviation to pair with the trough and V1 refuses with
\texttt{inadmissible\_precedence}. The figure
reproduces what the artefact contains; it is not a depiction of the
incident's traffic impact.}
\label{fig:case}
\end{figure}

\subsection{What the system says, and what that is worth}

Under the Phase-1 mapping, a window with routing detections and no
traffic observation is the \texttt{silent\_alert} pattern: something
moved in the control plane and the data plane did not register it. For
this event that pattern is the correct \emph{description} of the
telemetry, and it is also the correct description of the incident, for
different reasons---the traffic series has no baseline, and the
incident's data-plane effect is a validation-driven unreachability
rather than a volume change at the vantage point. A verdict that
conflates the two would be right for the wrong reason, which is the
failure mode \S\ref{sec:intro} attributes to single-signal systems.

Phase~2 abstains, and \S\ref{sec:eval-shape} gives the arithmetic. The
50-entry timeline becomes a grid of 91 buckets (the engine's own
five-minute bucketing) or 120 (a 30-second grid, capped), with exactly
two non-filler buckets in either case: the routing spike at one end, the
synthesized traffic record at the other---a point placed at the midpoint
of a 33-hour traffic window, 7.5 hours from the only routing evidence.
CSS returns $p = 1.0$ rather than attribute a traffic pre/post shift of
$0.11\sigma$; transfer entropy estimates $0.03$\,nats in the $R\to T$
direction against a shuffle null of $0.061\pm0.004$ ($p = 1.0$). The
engine does not reach either statistic: with one non-filler point per
side there is no sustained deviation to pair with the trough, so V1 returns
\texttt{inadmissible\_precedence} with sub-verdict
\texttt{inadmissible\_precedence}.

It is worth being explicit about why that abstention is the right output
rather than a missed opportunity. The ground truth for this event contains
no traffic anomaly. An engine that committed $R\to T$ here because the
point estimates happened to fall that way would have been \emph{right by
luck}: right about the incident, wrong about the evidence, and therefore
wrong on the next event with the same shape but the opposite cause.

\subsection{Lessons}

Three artefacts in one window masqueraded as evidence: a detector firing
without its baseline (\texttt{expected\_origin} = null), a traffic
verdict computed against a baseline that does not exist
(\texttt{historical\_avg} = 0.0), and a synthesized record standing in
for an observation (\texttt{cluster\_anomaly\_count} = 0 with a non-zero
severity). Each is individually plausible and each would have survived a
review that looked only at the verdict. The reusable lesson is that a
refusal is only as good as the provenance it records: the vetoes can
decline to commit, but they cannot know that a firing detector had no
baseline unless the cache says so. That is why the three-state recording
of \S\ref{sec:impl} matters more than any single accuracy number in this
paper, and why the case study is a study of instrumentation rather than
of attribution.
